\section{Discussion}
\label{sec:discussion}

\subsection{Interpretation of Quality Results}

The 29\% relative quality improvement of the full pipeline over the No-Refs ablation
demonstrates that reference verification — not prose fluency or visual richness — is
the dominant quality driver in automated academic document generation.
This finding has a principled explanation in the structure of the quality metric $Q$
(Equation~\ref{eq:quality}): the reference integrity sub-metric $R(s)$ (weight $w_2=0.35$)
is exactly the dimension that the \cgrag{} gate is designed to maximise, and the
gate's effect propagates downstream because $F_3$ (drafting) uses the verified reference
set $\mathcal{R}^+$ as its primary context source.

The propagation mechanism is straightforward: when the drafting stage receives a
high-quality, verified context window, it generates prose that cites the verified
references explicitly and consistently.
Sentence-level citation patterns learned from the training corpus of the LLM cause
it to insert \texttt{\textbackslash cite\{...\}} commands that reference entries
in $\mathcal{R}^+$.
Because $\mathcal{R}^+$ contains only \texttt{enhanced} references, the resulting
citations are all resolvable at validation time, achieving near-ceiling $R(s_5)$ scores.

In the No-Refs ablation, the opposite dynamic operates.
Without $F_2$, the drafting stage has no verified reference context.
The LLM generates plausible-sounding citations by pattern-matching to scholarly
citation formats it has seen during training, but many of these fabricated citations
do not correspond to real publications.
The result is high surface-level fluency (the prose reads naturally) combined with
low $R(s)$ (most citations fail the live-resolve check).
This is precisely the hallucination failure mode identified in the literature~\cite{alkaissi2023hallucinated},
and the \cgrag{} gate is the mechanism that blocks it.

The 10\% improvement of the full pipeline over the No-Visuals ablation is smaller
but consistent and attributable to a different mechanism.
When figures are absent, body text that contains \texttt{\textbackslash ref\{fig:*\}}
directives produces unresolved references at compile time.
These unresolved references cause a visible ``??'' in the compiled PDF at the
reference location, penalising the visual coherence sub-metric $V(s)$.
The \tvg{} cycle prevents this: by generating compilable figures with matching
\texttt{\textbackslash label\{fig:*\}} entries before assembly, it ensures that
all figure references in the body resolve to actual figure environments.

\subsection{Why the Formal Model Matters}

One might ask whether the formal specification of \sys{} is merely descriptive —
a post-hoc rationalisation of engineering decisions already made — or whether it
is prescriptive in the sense of enabling design decisions that could not have been
made without the formal framework.

We argue that the formal model is prescriptive in at least three respects.

First, the \textbf{pipeline algebra} $\Pi$ (Definition~\ref{def:pipeline}) prescribes
the type system for inter-stage data.
Without a formal type for each stage's output, the practical engineering question
of ``what does Stage~$F_3$ receive as input?'' is answered informally (e.g., by a
shared README or a verbal convention among developers).
Informal answers are fragile: they break when developers change the format of a stage's
output without updating the informal documentation.
The formal type constraint $\mathrm{cod}(F_i)=\mathrm{dom}(F_{i+1})$ makes this
break detectable at the specification level before any code is written.

Second, the \textbf{\cgrag{} gate} (Definition~\ref{def:cgrag}) prescribes the exact
conjunction of admission criteria.
An engineering team without a formal model might implement reference filtering as a
single threshold on confidence score, or as a weighted sum of criteria.
The lattice formulation reveals that a weighted sum is inadequate because it permits
compensatory admission (a very confident reference with zero citations is admitted),
which violates the design intent.
The meet-semilattice structure forces the designer to choose a conjunction, which
is the only algebraically correct implementation of ``all criteria must be satisfied.''

Third, the \textbf{billing operator} (Definition~\ref{def:billing}) prescribes the
linearity constraint.
An engineering team implementing billing without a formal model might use a piecewise
billing function (e.g., a lower rate for the first 5,000 tokens, a higher rate above
that threshold).
Invariant~I3 (Billing Monotonicity) immediately reveals that a piecewise function with
a rate change at 5,000 tokens is non-monotone around that threshold: spending exactly
5,000 tokens and then 1 more would result in a higher cost than if the boundary
were crossed differently.
The formal constraint forces the designer to choose a globally monotone function.

\subsection{Scope and Limitations of the Formal Model}

The Unified System Model $\mathcal{P}$ deliberately abstracts over several aspects
of the physical system.
We discuss each abstraction and its implications.

\paragraph{LLM non-determinism.}
The pipeline algebra $\Pi$ treats each stage $F_i$ as a deterministic morphism in
$\mathbf{DocState}$.
In practice, LLM calls are stochastic: the same prompt can produce different outputs
on different calls due to temperature-based sampling.
The formal model is valid in the sense that $\Pi$ specifies the \emph{intended}
transformation; the stochastic LLM is treated as a black-box realisation of that
transformation.
This is the same abstraction used in randomised algorithm analysis, where a
randomised algorithm is modelled as a deterministic algorithm applied to a random
tape.
Proposition~\ref{prop:mdp} provides the probabilistic extension that makes the
stochasticity explicit at the MDP level.

\paragraph{Network liveness.}
Invariant~I5 (Real-Data Guarantee) is conditional on network availability.
The formal statement $\forall\,r\in\mathrm{refs}(d)$: $r$ is live at $t_0$
is logically equivalent to: (network is live at $t_0$) $\Rightarrow$ ($r$ was fetched
at $t_0$).
The two I5 violations in the corpus (Section~\ref{sec:results}) are consistent with
the antecedent being false (the CDN cache violated the freshness protocol, so the
network was not ``live'' in the required sense).

A production-grade formalisation would introduce an explicit liveness predicate
$\lambda(t_0)\in\{\top,\bot\}$ and restate I5 as:
$\lambda(t_0)\Rightarrow\forall\,r\in\mathrm{refs}(d): r$ was fetched at $t_0$.
This makes the conditional nature of the invariant explicit and admits formal
verification against a TLA+ model of the network protocol.

\paragraph{LLM API evolution.}
The billing coefficients $(3, 800)$ in $\mathcal{B}$ are calibrated to the
\texttt{gpt-5-mini} API version deployed during the evaluation period.
A provider price change would require updating the coefficients without affecting
the formal structure of $\mathcal{B}$ (which requires only that the function is
affine with non-negative coefficients).
The formal model is therefore stable under price changes, which is a desirable
property: the specification does not need to be rewritten when the API changes.

\paragraph{Semantic correctness.}
None of the five invariants and neither of the two theorems make claims about the
\emph{semantic} correctness of the generated prose — only about structural properties
(compilability, citation resolvability, order, cost monotonicity).
Semantic correctness (whether the generated claims are factually accurate)
is outside the scope of formal verification at the current state of the art.
The \cgrag{} gate provides a proxy for semantic quality by filtering on scholarly
authority, but scholarly authority is not identical to factual accuracy:
a highly cited paper can contain incorrect claims.
We acknowledge this limitation and leave semantic verification to future work.

\subsection{The \cgrag{} Gate: Threshold Sensitivity}

The three thresholds $(\tau_c,\tau_s,\tau_k)=(7.0,100,3)$ were set based on a pilot
study of 40 sessions conducted prior to the main evaluation.
We report sensitivity analysis results for $\tau_c$ (the most influential threshold).

A grid search over $\tau_c\in\{5.0, 6.0, 7.0, 8.0, 9.0\}$ with
$(\tau_s,\tau_k)=(100,3)$ fixed yields the following quality scores:

\begin{center}
\begin{tabular}{lrrr}
\toprule
$\tau_c$ & \% references admitted & $\bar{Q}(s_5)$ & $\bar{R}(s_5)$ \\
\midrule
5.0 & 74\% & 0.93 & 0.88 \\
6.0 & 61\% & 0.96 & 0.93 \\
\textbf{7.0} & \textbf{48\%} & \textbf{0.97} & \textbf{0.96} \\
8.0 & 32\% & 0.94 & 0.97 \\
9.0 & 18\% & 0.88 & 0.98 \\
\bottomrule
\end{tabular}
\end{center}

The quality curve is unimodal with a maximum at $\tau_c=7.0$.
Lowering $\tau_c$ to $5.0$ admits $74\%$ of candidate references, increasing
context richness but introducing low-confidence references that reduce $R(s_5)$
from $0.96$ to $0.88$.
Raising $\tau_c$ to $9.0$ admits only $18\%$ of references — a context-impoverishment
effect that forces the drafting stage to work with insufficient source material,
reducing $\bar{Q}(s_5)$ to $0.88$ despite near-perfect reference integrity ($\bar{R}=0.98$).
The sensitivity analysis confirms that the current threshold $\tau_c=7.0$ is at the
optimum of the quality-integrity trade-off curve.

\subsection{Comparison with Embedding-Based Reference Selection}

To assess the design choice of using $\Gamma$ (multi-criteria lattice) rather than
embedding-cosine retrieval as in standard RAG, we ran a post-hoc comparison on the
same 120 sessions.
We replaced $F_2$ with a standard dense retrieval step using cosine similarity to
retrieve the top-10 references from a pre-indexed corpus of 2 million Semantic Scholar
abstracts, without citation count or keyword filtering.
The result was $\bar{Q}(s_5)=0.84$ — a 13-percentage-point reduction compared to the
full pipeline and a 9-percentage-point reduction compared to the No-Refs ablation.

The primary source of degradation was citation integrity: $\bar{R}(s_5)$ dropped
from $0.96$ (full pipeline) to $0.72$ (embedding-based).
Manual inspection of 20 randomly selected sessions from the embedding-based condition
revealed that approximately 35\% of admitted references were informal web pages
(Wikipedia articles, blog posts, documentation pages) that achieved high embedding
similarity to the academic query but did not satisfy the scholarly authority criteria
of the \cgrag{} gate.
This validates the design of $\Gamma$ as a hard admission gate rather than a soft
ranking criterion: embedding similarity is a necessary but not sufficient condition
for scholarly reference quality.

\subsection{Language Partition and Multilingual Quality}

Sessions assigned to $L_{\mathrm{lat}}$ (Kazakh Latin, $n=5$) achieved the lowest
$\bar{Q}(s_5)=0.91$ and $\bar{L}=0.82$ in the corpus.
Manual analysis of the 5 sessions reveals two distinct failure modes.

The first failure mode is \emph{script inconsistency}: the LLM generates body text
that switches between Kazakh Latin and Russian Cyrillic for technical terms.
This is attributable to the relative scarcity of Kazakh Latin text in the LLM's
training corpus: the 2021 Kazakh Latin alphabet was only recently standardised,
and pre-2021 training data contains essentially no Kazakh Latin text.
The language detection step correctly identifies $L_{\mathrm{lat}}$ at the start of
$F_1$, but the prompting directive (``respond in Kazakh Latin script'') is not
always followed for technical terminology.

The second failure mode is \emph{transliteration inconsistency}: when the LLM does
use Kazakh Latin, it sometimes applies the pre-2021 transliteration convention
(using apostrophes for palatalisation, e.g., \textit{j'}\ for the soft j sound)
instead of the 2021 standard (using special letters like \textit{\'j}).
This inconsistency does not affect compilability but does affect the $L(s)$ sub-metric.

Both failure modes can be addressed by a post-processing step that applies a
deterministic script-normalisation function to the draft output before assembly.
This is a planned improvement for a future version of the platform.

\subsection{Future Directions}

The formal model $\mathcal{P}$ opens four concrete directions for future work.

\begin{enumerate}[label=\textbf{F\arabic*.},leftmargin=2.4em]
  \item \textbf{Learned repair policy.}
    Proposition~\ref{prop:mdp} frames the repair cycle as an MDP.
    Training a learned policy that conditions on the semantic structure of the
    compiler error message (e.g., ``undefined colour name'' versus ``missing
    \texttt{\textbackslash end}'' versus ``undefined control sequence'') could
    provide a repair directive that is more targeted than the full stderr message,
    increasing $\hat{p}$ from $0.72$ to a higher value.
    A simple version of this policy could be a lookup table mapping error class
    to repair template, avoiding the need for a reinforcement-learning setup.

  \item \textbf{Adaptive $\tau$ selection.}
    The \cgrag{} thresholds $(\tau_c,\tau_s,\tau_k)$ are currently fixed across
    all sessions.
    A Bayesian update rule that starts with the current thresholds as a prior
    and updates them based on the observed reference quality distribution of prior
    sessions in the same topic category could improve per-domain performance.
    For example, biomedical sessions consistently admit preprint-level references
    at $Q_3$ that perform well in practice; a domain-specific lower $\tau_s$
    for this category would increase context richness without degrading $R(s_5)$.

  \item \textbf{Formal verification of I5.}
    A TLA+ specification of the retrieval protocol — including the explicit liveness
    predicate $\lambda(t_0)$ and the \texttt{If-Modified-Since} header check —
    could provide machine-verified proof of Invariant~I5 under an explicit network
    liveness assumption.
    This would replace the informal proof sketch in Section~\ref{sec:architecture}
    with a formally verified certificate, suitable for inclusion in a system safety
    argument.

  \item \textbf{Cross-lingual LCS reconciliation.}
    The current $\Phi_\theta$ operates on token sequences from the same language.
    Extending it to cross-lingual operation (aligning a Kazakh reference sentence
    with an English draft sentence) requires a cross-lingual token alignment that
    preserves the order-sensitivity property of LCS.
    One approach is to project both sequences into a shared multilingual token
    space using a multilingual encoder, then apply LCS in the projected space.
    The threshold $\theta$ would need to be recalibrated for the projected space,
    since projected token matches may have different semantic weight than surface-form
    token matches.
\end{enumerate}
